Puji syukur penulis panjatkan ke hadirat Allah SWT atas segala rahmat,
nikmat, dan hidayah-Nya sehingga penulis dapat menyelesaikan tugas akhir
berupa skripsi ini dengan baik dan lancar. Dalam perjalanan akademis
hingga penyusunan tugas akhir ini, penulis memperoleh banyak bimbingan,
bantuan, dukungan, serta motivasi dari berbagai pihak. Oleh karena itu,
pada kesempatan ini penulis ingin menyampaikan ucapan terima kasih kepada:

\begin{enumerate}

  \item Bapak Prof.\ Ir.\ Hanung Adi Nugroho, S.T., M.Eng., Ph.D., IPM.,
        SMIEEE.\ selaku Ketua Departemen Teknik Elektro dan Teknologi
        Informasi Fakultas Teknik Universitas Gadjah Mada.

  \item Bapak Dr.\ Ahmad Nasikun, S.T., M.Sc.\ selaku Ketua Program Sarjana
        Program Studi Teknologi Informasi Fakultas Teknik Universitas Gadjah
        Mada.

  \item Bapak Dani Adhipta, S.Si., M.T.\ selaku dosen pembimbing pertama
        yang telah memberikan banyak bantuan, bimbingan, dukungan, serta
        arahan kepada penulis selama proses penyusunan skripsi ini.

  \item Bapak Ir.\ Sujoko Sumaryono, M.T.\ yang telah memberikan arahan,
        masukan, dan bimbingan kepada penulis pada tahap awal penyusunan
        skripsi ini.

  \item Kedua orang tua tercinta, Ibu Sri Dayanti, S.Si.\ dan Bapak Dr.\
        Ir.\ Arif Wismadi, M.Sc.\ yang jasanya tidak akan pernah dapat
        terbalaskan serta senantiasa memberikan doa, dukungan, kasih sayang,
        dan motivasi terbaik bagi penulis. Serta Fathin Sabbiha Wismadi, S.Ars., M.Sc., M.\ Tsaqif Wismadi, S.PWK., M.Sc., Ph.D.,
        dan Hanifah Nabila Wismadi, S.T.\ yang selalu membersamai dan
        memberikan dukungan kepada penulis sepanjang proses hidup hingga
        penyusunan skripsi ini.

  \item Yitzhak, Grandiv, Sakti, Bian, dan Yahya serta seluruh tim
        ArachnoVa lain yang telah menjadi rekan berbisnis, berdiskusi,
        bekerja sama, serta memberikan pengalaman dan dukungan selama
        perjalanan akademis penulis.

  \item Luthfi dan Varick yang telah membersamai penulis dalam berbagai
        proses akademik, kerja praktik, serta menjadi teman sejak awal
        perkuliahan hingga proses penyusunan skripsi ini.

  \item Tim Capstone E-06, yaitu Bayu, Fata, Charlenne, dan Malika yang
        telah menjadi rekan seperjuangan dalam menyelesaikan berbagai
        rangkaian kegiatan akademik pada program Capstone C251 dan C501.

  \item Tim KKN Kisah Panorama 2025, yaitu Pram, Farros, Alan, Izzat,
        dan Abi, serta teman-teman lainnya yang tidak dapat disebutkan
        satu per satu, yang telah bersama-sama menjalani dan membersamai
        kegiatan Kuliah Kerja Nyata dengan penuh kebersamaan dan
        pengalaman berharga.

  \item Teman-teman WEES yang memberikan kebersamaan, dan semangat kepada
        penulis selama masa perkuliahan hingga penyusunan skripsi ini.

  \item Seluruh dosen dan tenaga kependidikan DTETI FT UGM yang telah
        memberikan ilmu, pengalaman, serta membantu kelancaran penulis
        sejak awal hingga akhir masa perkuliahan.

  \item Seluruh teman-teman seperjuangan di DTETI UGM serta berbagai pihak
        lainnya yang tidak dapat disebutkan satu per satu. Terima kasih atas
        kebersamaan, bantuan, serta kenangan yang telah diberikan kepada
        penulis.

\end{enumerate}

Akhir kata, penulis berharap semoga skripsi ini dapat memberikan manfaat
bagi semua pihak yang membacanya.

\begin{flushright}
\vspace{\baselineskip}
Yogyakarta, 1 Juni 2026

\vspace{\baselineskip}
Penulis
\end{flushright}
